\chapter{\uppercase{Background and Related Work}}
% Page numbering handled by Chapter 1
\doublespacing

% ============================================================================
% CHAPTER 2: BACKGROUND AND RELATED WORK
% Scope: GNNs for vulnerability detection, LLMs for code security,
%         conformal prediction, multi-agent systems, hybrid cascaded approaches.
% NOTE: Pattern-based SAST, taint analysis, CodeQL, and CPGs were covered
%       in Phase 1 Chapter 2 and are NOT duplicated here.
% ============================================================================

The preceding chapter established that static analysis tools suffer from
false positive rates exceeding 90\% at scale~\cite{ghostsecurity2025},
and that the resulting alert fatigue causes nearly half of all discovered
vulnerabilities to remain unpatched for over twelve
months~\cite{johnson2013}. We formulated three research questions
targeting uncertainty-driven escalation, conformal coverage guarantees,
and adversarial dual-agent validation. To ground these questions in the
state of the art, this chapter surveys five areas that directly inform
our framework design: graph neural networks for vulnerability detection
(\S\ref{sec:gnn-lit}), large language models for code security
(\S\ref{sec:llm-lit}), conformal prediction and calibrated uncertainty
(\S\ref{sec:conformal-lit}), hybrid and cascaded detection systems
(\S\ref{sec:hybrid-lit}), and the research gaps that our contributions
address (\S\ref{sec:gaps}).

% ============================================================================
\section{\uppercase{Graph Neural Networks for Vulnerability Detection}}
\label{sec:gnn-lit}
% ============================================================================

\subsection{From Sequential Tokens to Graph Representations}

Traditional machine-learning approaches to vulnerability detection treat
source code as a flat token sequence and apply convolutional or recurrent
models~\cite{thapa2022}. This representation discards the structural
relationships---control flow, data dependence, function calls---that
determine whether a code pattern is actually exploitable. Code Property
Graphs (CPGs), introduced by Yamaguchi et al.~\cite{cpg2014} and
surveyed in our Phase~1 report, unify Abstract Syntax Trees, Control
Flow Graphs, and Program Dependence Graphs into a single heterogeneous
graph. Graph neural networks (GNNs) operate directly on these graph
structures, propagating information along edges to learn representations
that encode both syntactic form and semantic relationships.

\subsection{Foundational GNN Architectures}

\textbf{Devign}~\cite{devign2019} was among the first to apply GNNs to
vulnerability detection at scale. Zhou et al.\ constructed joint graphs
combining AST, CFG, DFG, and Natural Code Sequence (NCS) edges, then
trained a Gated Graph Neural Network (GGNN) on function-level
classification. On a dataset of 27,318 functions from QEMU and FFmpeg,
Devign achieved 63.69\% accuracy. While this result established the
viability of graph-based detection, the absolute performance left
substantial room for improvement, and the reliance on GGNN---which uses
a fixed number of message-passing iterations---limited the model's
ability to capture long-range dependencies in large functions.

\textbf{Reveal}~\cite{reveal2021} (Chakraborty et al., IEEE TSE 2021)
advanced the field by directly confronting the class imbalance problem
that plagues vulnerability datasets. Real-world code is overwhelmingly
non-vulnerable; na\"{\i}ve training produces classifiers biased toward
the majority class. Reveal proposed a CPGNN architecture with
triplet-loss training and data-rebalancing strategies, demonstrating that
dataset quality matters as much as model architecture. Critically,
Chakraborty et al.\ showed that existing datasets contained systematic
label noise, an observation later amplified by the PrimeVul
study~\cite{primevul2024}.

\textbf{IVDetect}~\cite{ivdetect2021} (Li et al., ESEC/FSE 2021)
introduced interpretation-based vulnerability detection with program
slicing. Rather than classifying an entire function, IVDetect identifies
the specific statements contributing to a vulnerability through GNN
attention weights. This line-level interpretability is valuable for
developer triage, as it narrows the scope of manual review. We adopt a
similar principle in our framework: the MiniGINv3 attention scores over
CPG nodes provide localization cues that are propagated to the final
SARIF report.

\textbf{ReGVD}~\cite{regvd2022} (Nguyen et al., IEEE TSE 2022) revisited
GNN design choices for vulnerability detection, evaluating multiple
graph construction strategies and GNN backbones. Their key finding was
that Graph Isomorphism Networks (GIN) with code-gadget representations
achieved competitive F1 scores ($\sim$0.62--0.68 on Devign). This
result is theoretically grounded: Xu et al.~\cite{xu2019gin} proved that
GIN, with its injective sum aggregation, is as expressive as the
Weisfeiler-Lehman (WL) graph isomorphism test---the theoretical ceiling
for message-passing GNNs. In contrast, GAT uses attention-weighted mean
aggregation, which is not injective and can conflate structurally
distinct neighborhoods. We leverage this theoretical insight in our
MiniGINv3 design: we adopt GIN's injective sum aggregation for its
superior expressiveness, combined with attention-based readout to
retain the localization capability required for vulnerability
triage.

\subsection{Transformer-Based Alternatives}

\textbf{LineVul}~\cite{linevul2022} (Fu and Tantithamthavorn, MSR 2022)
took a different approach, applying a RoBERTa-based transformer to
vulnerability prediction at line-level granularity. Operating on the
BigVul dataset, LineVul achieved strong reported F1 scores and
demonstrated that pre-trained language models can capture vulnerability
patterns without explicit graph construction. The key advantage of
LineVul is simplicity: it requires no CPG extraction, making it
applicable to any language with a tokenizer. The key disadvantage is
that it treats code as a linear token sequence, discarding the structural
information that distinguishes, for example, a sanitized input flow from
an unsanitized one.

\textbf{GraphCodeBERT}~\cite{graphcodebert2021} (Guo et al., ICLR 2021)
bridges the gap between sequential and structural representations by
incorporating data-flow information into the pre-training objective.
Unlike CodeBERT~\cite{codebert2020}, which treats code as natural
language, GraphCodeBERT learns edge-level data-flow relationships during
pre-training, producing embeddings that are structurally aware without
requiring a full CPG at inference time. We use GraphCodeBERT as the
feature extractor in our graph pipeline: each code node receives a
768-dimensional embedding from GraphCodeBERT, concatenated with 6
structural features (normalized in-degree, normalized out-degree, sink indicator,
source indicator, normalized BFS depth, and language identifier), yielding 774-dimensional
input vectors for MiniGINv3.

\subsection{The Benchmark Inflation Problem}

The most consequential recent finding in this area comes from
\textbf{PrimeVul}~\cite{primevul2024} (Ding et al., ICSE 2025), which
demonstrated that reported GNN and transformer performance on standard
benchmarks is dramatically inflated. A 7B-parameter model that scored
68.26\% F1 on the BigVul benchmark achieved only 3.09\% F1 on
PrimeVul's rigorously deduplicated C/C++ dataset. The collapse is caused
by data leakage: near-duplicate functions appear in both training and
test sets, and models memorize surface patterns rather than learning
genuine vulnerability semantics. This finding has two implications for
our work. First, we cannot rely on benchmark F1 as the sole evaluation
metric; any claimed detection performance must be validated on
deduplicated data with temporal splits. Second, the fragility of
point-prediction models motivates our use of conformal prediction
(\S\ref{sec:conformal-lit}), which provides coverage guarantees
independent of the model's raw accuracy.

\subsection{Comparative Analysis}

Table~\ref{tab:gnn-lit} summarizes the GNN-based approaches discussed
above. The consistent theme across these works is a trade-off between
representation richness and practical deployment. Richer graph
representations (full CPGs) improve detection but require expensive
extraction tooling. Simpler representations (token sequences, partial
ASTs) are more portable but miss structural patterns. None of these
systems provide calibrated uncertainty quantification---they output point
predictions without indicating when the model is unsure.

\begin{table}[H]
\centering
\caption{Comparison of GNN-based vulnerability detection approaches.}
\label{tab:gnn-lit}
\small
\renewcommand{\arraystretch}{1.3}
\begin{tabular}{p{1.8cm} p{2.0cm} p{1.8cm} p{1.5cm} p{2.0cm} p{3.2cm}}
\toprule
\textbf{Approach} & \textbf{Architecture} & \textbf{Dataset} & \textbf{Metric} & \textbf{Representation} & \textbf{Key Limitation} \\
\midrule
Devign \cite{devign2019} & GGNN & QEMU/FFmpeg & 63.7\% Acc & AST+CFG+DFG +NCS & Fixed iteration count; no uncertainty \\
Reveal \cite{reveal2021} & CPGNN + triplet loss & Reveal dataset & Improved over baselines & Full CPG & Requires Joern; no confidence score \\
IVDetect \cite{ivdetect2021} & GNN + slicing & BigVul & F1 reported & PDG slices & Expensive slicing; single language \\
ReGVD \cite{regvd2022} & GIN on gadgets & Devign & $\sim$0.62--0.68 F1 & Code gadgets & No inter-procedural flow \\
LineVul \cite{linevul2022} & RoBERTa & BigVul & Line-level F1 & Token sequence & No structural information \\
GraphCodeBERT \cite{graphcodebert2021} & Transformer + DFG & CodeSearchNet & Improved over CodeBERT & Tokens + data flow & Pre-training only; no CPG \\
PrimeVul \cite{primevul2024} & Evaluation study & PrimeVul & 3.09\% F1 (realistic) & Multiple & Evaluation, not a detection method \\
\bottomrule
\end{tabular}
\end{table}

% ============================================================================
\section{\uppercase{Large Language Models for Code Security}}
\label{sec:llm-lit}
% ============================================================================

\subsection{LLMs as Vulnerability Detectors}

The application of large language models to vulnerability detection
follows two paradigms: \emph{classification}, where the LLM labels code
as vulnerable or safe, and \emph{reasoning}, where the LLM explains
whether a vulnerability is exploitable. The classification paradigm
inherits the same benchmark-inflation problem identified by
PrimeVul~\cite{primevul2024}: GPT-3.5 and GPT-4 performed at levels
``akin to random guessing'' on PrimeVul's most stringent evaluation
setting. The reasoning paradigm is more promising but introduces its own
failure modes.

\textbf{LLM4Vuln}~\cite{llm4vuln2024} (Sun et al., 2024) proposed a
unified evaluation framework that decouples vulnerability reasoning into
sub-tasks: detection, classification, and explanation. By structuring
the LLM interaction as a multi-role pipeline, LLM4Vuln achieved improved
consistency over single-prompt approaches. The framework demonstrated
that breaking complex security analysis into focused sub-tasks reduces
hallucination and improves the quality of generated explanations. We
adopt a similar decomposition principle: our Attacker and Defender agents
each focus on a specific analytical perspective rather than attempting
holistic assessment in a single prompt.

\textbf{Steenhoek et al.}~\cite{steenhoek2024} (2024) conducted a
systematic study of LLM capabilities and limitations for vulnerability
detection. Their findings confirmed that LLMs exhibit high
non-determinism---producing inconsistent classifications across repeated
runs---and that performance degrades on code patterns absent from
training data. These results underline the need for consensus mechanisms
that stabilize LLM outputs, which is precisely the function of our
dual-agent protocol.

\textbf{VulnHuntr}~\cite{vulnhuntr2024} (Protectai, 2024) demonstrated
that LLM-driven autonomous vulnerability hunting can discover real
zero-day vulnerabilities in open-source Python projects. VulnHuntr uses
iterative code traversal guided by LLM reasoning to trace data flows
across files and identify injection, deserialization, and SSRF
vulnerabilities. The system's success validates the LLM-as-reasoning-engine
approach, but its lack of structured validation means that reported
findings still require manual confirmation. In contrast, our framework
provides structured evidence from both the Attacker and Defender agents,
producing an auditable decision trail.

\subsection{Benchmark Fragility}

Thapa et al.~\cite{thapa2022} (ACSAC 2022) compared transformer-based
models (BERT, RoBERTa, GPT-2) on vulnerability detection, reporting
competitive results on standard datasets. However, the PrimeVul
study~\cite{primevul2024} subsequently demonstrated that these results
do not transfer to realistic data. The gap between benchmark and
real-world performance---from 68\% F1 to 3\% F1---represents the single
most important finding in recent vulnerability detection research. Any
system that relies exclusively on a single model's classification output,
whether GNN or LLM, inherits this fragility. Our cascade architecture
mitigates the risk by combining three independent analysis modalities,
each with different failure modes, and by using conformal prediction to
quantify the GNN's uncertainty rather than trusting its point prediction.

\subsection{Multi-Agent LLM Systems}

Single-model LLM inference suffers from confirmation bias: once the
model commits to an initial assessment, subsequent reasoning tends to
reinforce rather than challenge it. Multi-agent debate addresses this
limitation by introducing adversarial perspectives.

\textbf{Du et al.}~\cite{multiagentdebate2023} (2023) demonstrated that
multi-agent debate---where multiple LLM instances argue over an
answer---improves factual accuracy by 20--40\% over single-model
inference across mathematical reasoning, strategic analysis, and
factual question-answering tasks. The mechanism works because divergent
initial positions force each agent to generate counter-arguments,
exposing reasoning errors that a single model would propagate
unchallenged.

\textbf{Liang et al.}~\cite{divergentthinking2023} (2023) extended this
work by explicitly encouraging divergent thinking in multi-agent debate.
Their ``tit for tat'' protocol, where agents take turns challenging each
other's positions, improved performance on tasks requiring creative
reasoning. The key insight---that productive disagreement yields better
outcomes than artificial consensus---directly motivates our Attacker vs.\
Defender design, where the two agents are structurally incentivized to
reach opposing conclusions before a consensus engine reconciles them.

Our dual-agent protocol adapts these general multi-agent findings to
the specific structure of security analysis. The Attacker agent is
prompted with an offensive framing: given a code finding, construct a
concrete exploit scenario, identify the attack surface, and assess
exploitability. The Defender agent receives a defensive framing: identify
sanitizers, access controls, framework protections, and assess whether
the vulnerability is mitigated. This adversarial structure maps naturally
to the red team / blue team methodology that is standard practice in
security assessment, but---to our knowledge---has not previously been
applied to automated LLM-based vulnerability triage.

\subsection{Retrieval-Augmented Generation for Security}

LLMs possess broad but shallow knowledge of specific vulnerability
patterns. A model may know the general concept of SQL injection but lack
details about CVE-2023-specific bypass techniques or CWE-79 variant
classifications. Retrieval-Augmented Generation (RAG) addresses this by
injecting relevant knowledge at inference time.

Our framework maintains a knowledge base of over 200,000 NVD
entries~\cite{nvd2023} and 900+ CWE definitions, indexed using a hybrid
FAISS vector store and BM25 keyword retrieval with Reciprocal Rank
Fusion. When a finding escalates to the LLM stage, the RAG system
retrieves the most relevant CVE examples and CWE descriptions, providing
both agents with identical domain-specific context. This ensures that
the Attacker agent's exploit analysis and the Defender agent's
mitigation analysis are grounded in the same factual knowledge base
rather than depending on the LLM's parametric memory alone.

\subsection{Comparative Analysis}

Table~\ref{tab:llm-lit} summarizes the LLM-based approaches discussed
above. The central tension in this literature is between the semantic
reasoning power of LLMs and their unreliability as standalone
classifiers. LLMs excel at explaining \emph{why} code might be
vulnerable but fail at reliably determining \emph{whether} it is
vulnerable. Our framework exploits this asymmetry: we use the LLM as a
reasoning engine for the hardest cases (those that survived SAST and GNN
stages) rather than as a first-line classifier.

\begin{table}[H]
\centering
\caption{Comparison of LLM-based approaches for code security.}
\label{tab:llm-lit}
\small
\renewcommand{\arraystretch}{1.3}
\begin{tabular}{p{2.0cm} p{1.8cm} p{1.6cm} p{3.0cm} p{3.5cm}}
\toprule
\textbf{Approach} & \textbf{Model(s)} & \textbf{Dataset} & \textbf{Key Finding} & \textbf{Key Limitation} \\
\midrule
LLM4Vuln \cite{llm4vuln2024} & Multi-role LLMs & Custom & Decomposed reasoning improves consistency & No adversarial validation \\
PrimeVul \cite{primevul2024} & GPT-3.5/4, 7B models & PrimeVul & 68\% F1 $\to$ 3\% F1 on realistic data & Evaluation study, not a solution \\
VulnHuntr \cite{vulnhuntr2024} & LLM-driven & Python OSS & Real 0-days discovered & No structured validation; manual confirmation needed \\
Steenhoek et al.\ \cite{steenhoek2024} & Multiple LLMs & Multiple & 47--76\% inconsistency across runs & Study, not a detection method \\
Thapa et al.\ \cite{thapa2022} & BERT, RoBERTa, GPT-2 & Mixed & Competitive on benchmarks & Results do not transfer to realistic data \\
Du et al.\ \cite{multiagentdebate2023} & Multi-agent debate & Reasoning tasks & 20--40\% accuracy improvement & Not applied to code security \\
Liang et al.\ \cite{divergentthinking2023} & Divergent debate & Reasoning tasks & Divergent positions improve outcomes & Not applied to code security \\
\bottomrule
\end{tabular}
\end{table}

% ============================================================================
\section{\uppercase{Conformal Prediction and Calibrated Uncertainty}}
\label{sec:conformal-lit}
% ============================================================================

\subsection{The Calibration Problem in Neural Networks}

Modern deep neural networks are poorly calibrated: the probability
scores they output do not correspond to true likelihoods of correctness.
\textbf{Guo et al.}~\cite{guocalibration2017} (ICML 2017) demonstrated
this systematically, showing that as network depth and width increase,
calibration worsens---high-confidence predictions are frequently wrong.
In safety-critical applications such as vulnerability detection, this
miscalibration is dangerous: a model that reports 95\% confidence in a
``safe'' classification may be wrong 30\% of the time, leading to
undetected vulnerabilities in production code.

Post-hoc calibration methods such as temperature scaling~\cite{guocalibration2017}
and Platt scaling improve calibration on average but provide no
per-instance guarantees. A calibrated model may still produce arbitrarily
overconfident predictions on individual inputs that fall outside the
training distribution---precisely the ``novel code patterns'' that
represent zero-day vulnerabilities.

\subsection{Conformal Prediction: Distribution-Free Guarantees}

Conformal prediction, formalized by \textbf{Vovk et al.}~\cite{vovk2005}
(2005), provides a fundamentally different approach to uncertainty
quantification. Rather than producing a single prediction, conformal
prediction outputs a \emph{prediction set}---a subset of all possible
labels that is guaranteed to contain the true label with a user-specified
probability:

\begin{equation}
P\bigl(Y_{\text{true}} \in \mathcal{C}(X)\bigr) \geq 1 - \alpha
\label{eq:conformal_guarantee}
\end{equation}

\noindent where $\mathcal{C}(X)$ is the prediction set for input $X$
and $\alpha$ is the user-specified error rate. This guarantee is
\emph{distribution-free}: it holds for any data distribution, with the
only requirement being that the calibration data is exchangeable with the
test data. The guarantee holds in finite samples---no asymptotic
assumptions are needed.

\textbf{Angelopoulos and Bates}~\cite{angelopoulos2023} (2023) provided
a tutorial treatment that made conformal prediction accessible to the
machine learning community. They presented the split-conformal method,
where a held-out calibration set is used to compute conformity scores,
and a quantile threshold $\hat{q}$ is selected such that the empirical
coverage on the calibration set meets the target $1 - \alpha$. At test
time, the prediction set includes all labels whose conformity scores
exceed $\hat{q}$.

\subsection{Adaptive Prediction Sets}

The basic conformal prediction method can produce either very large sets
(uninformative) or very small sets (precise) depending on the score
function. \textbf{Romano et al.}~\cite{romano2020} (NeurIPS 2020)
introduced Adaptive Prediction Sets (APS), which use the model's own
softmax probabilities as conformity scores, sorted in decreasing order.
The prediction set is constructed by accumulating classes until the
cumulative probability exceeds $1 - \hat{q}$. APS produces smaller
prediction sets for inputs where the model is confident (sharp softmax
distributions) and larger sets where the model is uncertain (flat
distributions).

The APS method has a direct operational interpretation in our cascade.
For a binary vulnerability classifier:

\begin{itemize}
    \item A \textbf{singleton set} $\{$\texttt{safe}$\}$ or
          $\{$\texttt{vulnerable}$\}$ indicates that the GNN is
          confident enough to resolve the finding at Stage~2, because the
          conformal guarantee ensures the true label is included with
          probability $\geq 1 - \alpha$.
    \item A \textbf{two-element set} $\{$\texttt{safe},
          \texttt{vulnerable}$\}$ indicates genuine model uncertainty:
          both labels are plausible, and the finding must be escalated to
          Stage~3 (LLM dual-agent) for semantic resolution.
\end{itemize}

\noindent The singleton rate---the fraction of inputs receiving
single-label prediction sets---directly determines cascade efficiency.
A higher singleton rate means more findings are resolved at Stage~2
without expensive LLM calls.

\subsection{Conformal Prediction Under Distribution Shift}

A practical concern for deploying conformal prediction in vulnerability
detection is that the test distribution may differ from the calibration
distribution. New code patterns, frameworks, and vulnerability classes
appear continuously; the exchangeability assumption may be violated.
\textbf{Barber et al.}~\cite{barber2023} (Annals of Statistics, 2023)
extended conformal prediction beyond the exchangeability assumption,
providing coverage guarantees under distribution shift through
likelihood-ratio reweighting. Their results show that coverage can be
maintained even when the test distribution drifts from calibration data,
at the cost of wider prediction sets.

For our framework, this means that when the GNN encounters code patterns
dissimilar to its training data, the APS prediction set will
automatically widen (producing the ambiguous two-element set), triggering
escalation to the LLM stage. The conformal layer thus acts as an
automatic novelty detector: truly novel code is escalated by
construction, not by a hand-tuned threshold.

\subsection{Conformal Prediction in Security: A Gap}

We surveyed the conformal prediction literature across medical
imaging~\cite{angelopoulos2023}, autonomous driving, natural language
inference, and drug discovery. To our knowledge, no published work has
applied conformal prediction to vulnerability detection or static
analysis output triage. Existing vulnerability detection systems use
either hard classification thresholds (vulnerable if $p > 0.5$) or
ad-hoc confidence bins without formal coverage guarantees. Our
application of APS conformal prediction to the GNN output layer,
using singleton versus multi-label sets to drive cascade escalation,
is---to the best of our knowledge---the first such application in the
code security domain.

% ============================================================================
\section{\uppercase{Hybrid and Cascaded Detection Approaches}}
\label{sec:hybrid-lit}
% ============================================================================

\subsection{The Case for Multi-Stage Analysis}

The GNN and LLM limitations documented in \S\ref{sec:gnn-lit}
and~\S\ref{sec:llm-lit} motivate combining multiple analysis modalities.
No single technique---pattern matching, graph analysis, or LLM
reasoning---dominates across all vulnerability types. Pattern matching
excels at detecting known-bad API usage (hardcoded secrets, weak
cryptographic algorithms) but fails on context-dependent vulnerabilities.
Graph analysis captures structural properties (taint flow paths,
control-flow reachability) but struggles with semantic context
(is this sanitizer sufficient for this framework?). LLM reasoning
provides semantic understanding but at high cost and with unreliable
consistency. A multi-stage architecture that routes findings to the
appropriate analysis depth based on their characteristics can exploit
the complementary strengths of each modality.

\subsection{Existing Hybrid Systems}

\textbf{IRIS}~\cite{iris2025} (Li et al., ICLR 2025) augments CodeQL's
static analysis with LLM-inferred taint specifications. When CodeQL's
built-in models lack knowledge of a library's taint propagation behavior,
IRIS queries an LLM (GPT-4 or DeepSeek-Coder 7B) to infer whether
specific API methods are sources, sinks, or sanitizers. This approach
increased CodeQL's true positive detections by 103.7\% on
CWE-Bench-Java (120 real-world vulnerabilities). IRIS demonstrates that
LLMs can augment static analysis effectively, but it applies the LLM
\emph{uniformly} to all library specifications---there is no mechanism
to determine which queries warrant LLM augmentation and which do not.
Every finding incurs the same LLM cost regardless of analytical
difficulty.

\textbf{Vulnhalla} (CyberArk, 2025) post-filters CodeQL output through
an LLM-guided questioning protocol. For each CodeQL finding, the LLM
is presented with the flagged code and asked a structured series of
questions to determine whether the vulnerability is genuine. Vulnhalla
reported up to 96\% false positive reduction for specific CWE categories
on real GitHub repositories. The system validates our premise that LLM
reasoning can dramatically reduce SAST false positives, but it processes
\emph{every} finding identically through the LLM pipeline. For a scan
producing 1,000 findings, Vulnhalla makes 1,000 LLM calls---whereas our
cascade would make approximately 50, because 80\% of findings are
resolved at the SAST stage and a further 15\% at the GNN stage.

\textbf{ZeroFalse} (2025) enriches SAST trace data with contextual
information before submitting it to an LLM for adjudication. The system
reconstructs flow-sensitive traces from SAST output, adds surrounding
code context, and prompts a reasoning-oriented LLM to determine whether
each trace represents a true positive. ZeroFalse achieved 0.912 F1 on
the OWASP Java Benchmark and 0.955 F1 on an open-vulnerability dataset.
Like Vulnhalla, ZeroFalse sends all findings to the LLM without
selective escalation.

\textbf{LLMxCPG} (Lekssays et al., USENIX Security 2025) combines CPG
analysis with LLM classification in a two-phase architecture. In the
first phase, the LLM generates Joern queries to extract relevant CPG
slices. In the second phase, the LLM classifies the extracted slices as
vulnerable or safe. The approach achieved 15--40\% F1 improvement over
baselines on C/C++ code. LLMxCPG's use of CPG slicing to provide
structured context to the LLM validates our graph-then-LLM ordering,
but its two-phase design processes all findings through both phases
without cost-aware routing.

\subsection{Architectural Patterns}

Examining these systems reveals a consistent architectural pattern: all
existing hybrid approaches are \emph{two-stage} (SAST + LLM or
CPG + LLM) and \emph{uniform}---every finding receives the same
analysis regardless of its analytical difficulty. None incorporate a
GNN as an intermediate analysis layer. None use uncertainty scores to
drive escalation decisions. None provide formal coverage guarantees on
their predictions. Table~\ref{tab:gap-positioning} in
Section~\ref{sec:gaps} quantifies these differences.

% ============================================================================
\section{\uppercase{Research Gaps and Positioning}}
\label{sec:gaps}
% ============================================================================

The literature surveyed in the preceding sections reveals four
concrete gaps in the current state of the art. Each gap directly
motivates a corresponding contribution of our framework.

\subsection{Gap 1: No Selective Escalation Between Analysis Modalities}

Every hybrid system reviewed---IRIS, Vulnhalla, ZeroFalse,
LLMxCPG---applies its most expensive analysis stage uniformly to all
findings. IRIS queries the LLM for every unmodeled library.
Vulnhalla sends every CodeQL finding to its LLM pipeline. ZeroFalse
enriches every SAST trace for LLM adjudication. This uniform
application ignores a fundamental observation: the vast majority of
SAST findings are either clearly true positives (e.g., hardcoded
passwords) or clearly false positives (e.g., sanitized inputs) that do
not require expensive analysis. Only a minority of findings occupy an
ambiguous middle ground that benefits from deeper reasoning.

No published system uses a \emph{quantified uncertainty measure} to
determine which findings warrant escalation. Our \textbf{Contribution~1}
addresses this gap with a 4-factor uncertainty scoring model
($U = 0.4 \cdot U_{\text{conf}} + 0.3 \cdot U_{\text{comp}} + 0.2 \cdot
U_{\text{nov}} + 0.1 \cdot U_{\text{confl}} + \delta_{\text{sev}}$)
that routes findings through a 3-stage cascade. In our design, findings
with $U < 0.5$ are resolved at the SAST stage ($\sim$75--80\% of all
findings), reducing LLM API calls by approximately 85\% compared to
uniform approaches.

\subsection{Gap 2: No Calibrated Uncertainty in Vulnerability Detection}

All GNN-based and LLM-based vulnerability detectors reviewed in this
chapter produce point predictions---a single label or a probability
score---without formal calibration. Guo et al.~\cite{guocalibration2017}
showed that neural network confidence scores are systematically
miscalibrated, and PrimeVul~\cite{primevul2024} demonstrated that
vulnerability detection models are particularly fragile on
out-of-distribution data. A model reporting 90\% confidence that code
is safe may be wrong far more than 10\% of the time.

No vulnerability detection tool provides distribution-free coverage
guarantees. Our \textbf{Contribution~2} addresses this gap by applying
Adaptive Prediction Sets (APS)~\cite{romano2020} to the GNN output
layer, providing the guarantee in Equation~\ref{eq:conformal_guarantee}.
When the model is uncertain, the prediction set automatically expands,
triggering escalation. This replaces arbitrary confidence thresholds
with a mathematically principled decision boundary.

\subsection{Gap 3: LLM Triage Lacks Adversarial Validation}

Existing LLM-based triage systems (LLM4Vuln, VulnHuntr, Vulnhalla,
ZeroFalse) use a single LLM instance or a single-perspective prompt to
classify findings. Single-model inference is susceptible to confirmation
bias~\cite{steenhoek2024}: the model commits to an initial assessment
and reinforces it in subsequent reasoning. Multi-agent
debate~\cite{multiagentdebate2023, divergentthinking2023} has been shown
to improve LLM accuracy by 20--40\%, but this technique has not been
applied to code security triage.

Our \textbf{Contribution~3} addresses this gap with an adversarial
dual-agent protocol. The Attacker agent (red team) attempts to
construct exploits; the Defender agent (blue team) identifies
mitigations. A 5-rule consensus engine reconciles their perspectives
into a final verdict with structured evidence. This design maps
directly to the adversarial thinking (red team / blue team) that is
standard practice in manual security assessment but absent from
automated tools.

\subsection{Gap 4: Fixed Fusion Weights Ignore CWE-Specific Behavior}

Hybrid systems that combine outputs from multiple analysis stages
typically use fixed combination strategies: binary voting, simple
averaging, or a single set of learned weights. This approach ignores the
empirical observation that different vulnerability types respond
differently to different analysis methods. Injection vulnerabilities
(CWE-89, CWE-78, CWE-79) require semantic understanding of
sanitization context and are best served by LLM reasoning.
Cryptographic weaknesses (CWE-327, CWE-328) are identifiable by
syntactic patterns and are best served by SAST. Memory safety
vulnerabilities (CWE-416, CWE-476) require structural analysis of
allocation and deallocation paths and are best served by graph analysis.

No published system calibrates its fusion weights per CWE category.
Our \textbf{Contribution~4} addresses this gap with CWE-adaptive
score fusion, where the final confidence score
$f = \alpha \cdot S_{\text{SAST}} + \beta \cdot S_{\text{GIN}} +
\gamma \cdot S_{\text{LLM}}$ uses per-CWE calibrated weights
$(\alpha, \beta, \gamma)$ that reflect the empirical strengths of each
stage for that vulnerability class.

\subsection{Gap--Contribution Positioning}

Table~\ref{tab:gap-positioning} positions our framework against
existing systems across six capability dimensions. A checkmark
(\checkmark) indicates that the system provides the capability; a
dash (---) indicates absence.

\begin{table}[H]
\centering
\caption{Research gap positioning matrix. SEC-C is the only system providing all six capabilities.}
\label{tab:gap-positioning}
\small
\renewcommand{\arraystretch}{1.3}
\begin{tabular}{p{2.2cm} c c c c c c}
\toprule
\textbf{System} & \rotatebox{70}{\textbf{Uncertainty Routing}} & \rotatebox{70}{\textbf{Conformal Prediction}} & \rotatebox{70}{\textbf{Dual-Agent Validation}} & \rotatebox{70}{\textbf{CWE-Adaptive Fusion}} & \rotatebox{70}{\textbf{Multi-Language ($\geq$5)}} & \rotatebox{70}{\textbf{Free-Tier Operation}} \\
\midrule
Semgrep \cite{semgrep2024} & --- & --- & --- & --- & \checkmark & \checkmark \\
CodeQL \cite{codeql2024}   & --- & --- & --- & --- & \checkmark & \checkmark \\
Devign \cite{devign2019}   & --- & --- & --- & --- & --- & --- \\
LineVul \cite{linevul2022} & --- & --- & --- & --- & --- & --- \\
IRIS                       & --- & --- & --- & --- & --- & --- \\
Vulnhalla                  & --- & --- & --- & --- & --- & --- \\
ZeroFalse                  & --- & --- & --- & --- & --- & --- \\
LLMxCPG                    & --- & --- & --- & --- & --- & --- \\
\textbf{SEC-C (Ours)}      & \checkmark & \checkmark & \checkmark & \checkmark & \checkmark & \checkmark \\
\bottomrule
\end{tabular}
\end{table}

\noindent The positioning matrix makes the contribution space clear:
existing systems address subsets of the problem---SAST tools provide
rule-based detection, GNNs provide structural analysis, LLM systems
provide semantic reasoning---but no system integrates these modalities
with principled uncertainty-driven routing. SEC-C occupies a unique
point in this design space by combining all three analysis modalities
with conformal coverage guarantees, adversarial LLM validation, and
CWE-aware fusion into a single cascade architecture.

% ============================================================================
\section{\uppercase{Summary}}
\label{sec:litsurvey-summary}
% ============================================================================

This chapter surveyed the five research areas that inform our framework
design. GNN-based vulnerability detection
(\S\ref{sec:gnn-lit}) has produced architectures capable of learning
from graph-structured code representations, but reported performance is
severely inflated by benchmark artifacts~\cite{primevul2024}, and no GNN
system provides calibrated uncertainty. LLM-based code security
(\S\ref{sec:llm-lit}) brings semantic reasoning power but suffers from
inconsistency~\cite{steenhoek2024}, confirmation
bias~\cite{multiagentdebate2023}, and hallucination, making standalone
LLM classification unreliable. Conformal prediction
(\S\ref{sec:conformal-lit}) offers distribution-free coverage guarantees
that address the calibration gap, but it has not been applied to
vulnerability detection. Hybrid systems (\S\ref{sec:hybrid-lit}) combine
SAST with LLM reasoning to reduce false positives, but all existing
approaches process findings uniformly without cost-aware escalation.

We identified four research gaps: the absence of selective escalation
(Gap~1), the absence of calibrated uncertainty (Gap~2), the absence of
adversarial LLM validation (Gap~3), and the use of fixed fusion weights
that ignore CWE-specific behavior (Gap~4). These gaps map directly to
our four contributions: uncertainty-driven cascading escalation (C1),
APS conformal prediction for code security (C2), adversarial dual-agent
triage (C3), and CWE-adaptive score fusion (C4).

The next chapter presents the methodology of our framework, detailing
the architecture, algorithms, and design decisions that address each of
these identified gaps.

